\documentclass[9pt]{beamer}

\usepackage[english]{babel}
\usepackage[utf8]{inputenc}
\usepackage{hyperref}
\usepackage{verbatim}
\usepackage{multicol}
\usepackage{pifont}
\usepackage{wasysym}
\usepackage{color}

\definecolor{mygreen}{RGB}{27,156,16}
\definecolor{myred}{RGB}{221,0,17}
\newcommand{\lighter}[1]{{\color{gray} #1}}
\newcommand{\green}[1]{{\color{mygreen} #1}}
\newcommand{\red}[1]{{\color{myred} #1}}

\newcommand{\code}[1]{\texttt{#1}}
\newcommand{\outline}[0]{
  \frame{\frametitle{Outline}\tableofcontents[hideallsubsections]}
}
\newcommand{\outlinereminder}[0]{
  \frame{\frametitle{Outline}
  \tableofcontents[currentsection,subsectionstyle=show/show/hide]}
}



%      _         _             _                     _                     
%  ___| |_ _   _| | ___    ___| |_    ___ ___  _   _| | ___ _   _ _ __ ___ 
% / __| __| | | | |/ _ \  / _ \ __|  / __/ _ \| | | | |/ _ \ | | | '__/ __|
% \__ \ |_| |_| | |  __/ |  __/ |_  | (_| (_) | |_| | |  __/ |_| | |  \__ \
% |___/\__|\__, |_|\___|  \___|\__|  \___\___/ \__,_|_|\___|\__,_|_|  |___/
%          |___/                                                           
%

\setbeamercovered{transparent}
\usetheme{Warsaw}
% enlever [red] dans le \documentclass si utilisation de la suite :
\usecolortheme[RGB={202,31,39}]{structure} % rouge

%\setbeamertemplate{background}{\includegraphics[height=96mm,width=128mm]{fond.png}}

% definition de la couleur du fond des blocks
\definecolor{ArrierePlanBodyW}{rgb}{0.5,0.5,0.5} % si ecriture blanche
\definecolor{ArrierePlanBodyB}{rgb}{0.9,0.9,0.9} % si ecriture noire
% definition de la couleur des slides
\definecolor{ArrierePlan}{rgb}{1,1,1}

\definecolor{darkred}{rgb}{0.6,0,0}

\definecolor{toutnoir}{rgb}{0,0,0}

% selection de la couleur du fond des slides
\setbeamercolor{background canvas}{bg=ArrierePlan} 
% selection de la couleur du fond des block avec ecriture noire
\setbeamercolor{block body}{use=structure,fg=black,bg=ArrierePlanBodyB}
% selection de la couleur du fond des block avec ecriture blanche
%\setbeamercolor{block body}{use=structure,fg=white,bg=ArrierePlanBodyW}

%\setbeamercolor{section in head/foot}{toutnoir}

\setbeamertemplate{footline}{%
  \leavevmode%
  \hbox{\begin{beamercolorbox}[wd=.5\paperwidth,ht=2.5ex,dp=1.125ex,leftskip=.3cm plus1fill,rightskip=.3cm]{author in head/foot}%
    \usebeamerfont{author in head/foot}\insertshortauthor
  \end{beamercolorbox}%
  \begin{beamercolorbox}[wd=.5\paperwidth,ht=2.5ex,dp=1.125ex,leftskip=.3cm,rightskip=.3cm plus1fil]{title in head/foot}%
    \usebeamerfont{title in head/foot}\insertshorttitle \hfill \insertframenumber{} / \inserttotalframenumber\hspace*{2ex}%
  \end{beamercolorbox}}%
  \vskip0pt%
}


\setbeamercovered{dynamic}
%\logo{\includegraphics[height=1cm]{./img/logo_petit-tr.png}}
% \defbeamertemplate*{headline}{split theme}
% {%
%   \begin{beamercolorbox}{section in head/foot}
%     \insertsectionnavigationhorizontal{2cm}{\includegraphics[height=2cm]{./img/logo_petit-tr.png}}{}
%   \end{beamercolorbox}%
%  }

\defbeamertemplate*{headline}{}
{%
%   \begin{beamercolorbox}[wd=.1\paperwidth,ht=\tempdimb]{section in head/foot}
%      \vbox to 1cm{\vfil\includegraphics[.1\paperwidth]{./img/logo_petit-tr.png}\vfil}%
%   \end{beamercolorbox}
  \leavevmode%
  \begin{beamercolorbox}[wd=.095\paperwidth,ht=1.45cm]{section in head/foot}%
    \vbox to 1.45cm{\vfill\hspace*{1mm}\includegraphics[height=1.25cm]{lifc.png}\vfill}%
  \end{beamercolorbox}%
  \begin{beamercolorbox}[wd=.405\paperwidth,ht=1.45cm]{section in head/foot}%
    \vbox to 1.45cm{\vfil\insertsectionnavigation{.4\paperwidth}\vfil}%
  \end{beamercolorbox}%
  \begin{beamercolorbox}[wd=.33\paperwidth,ht=1.45cm]{subsection in head/foot}%
    \vbox to 1.45cm{\vfil\insertsubsectionnavigation{.4\paperwidth}\vfil}%
  \end{beamercolorbox}%
  \begin{beamercolorbox}[wd=.17\paperwidth,ht=1.45cm]{section in head/foot}%
%     \vbox to 1.5cm{\vfil\includegraphics[width=.145\paperwidth]{./img/ufc2.pdf}\vfill}%
%      \vbox to 1.5cm{\vfill\includegraphics[width=.185\paperwidth]{./img/ufc2.pdf}\vfill}%
      \vbox to 1.45cm{\vfill\hspace*{-0.1mm}\includegraphics[height=1.45cm]{ufc2.pdf}\vfill}%
  \end{beamercolorbox}%
}



\title{Grammar-Based Testing using \\ Realistic Domains in PHP}
\author{
    \textbf{Ivan~Enderlin},
    Frédéric~Dadeau,
    Alain~Giorgetti,
    Fabrice~Bouquet
}
\date{
    April 17th, 2012 \\
    A-MOST, Montréal
}


\begin{document}

\maketitle

\section{Introduction}

\begin{frame}
\frametitle{Context}

\begin{block}{Web}
\begin{itemize}
\item Its data and its languages: XML, HTML, forms, database queries,
network protocols…
\item Strings are the most used and manipulated data
\item They can be complex
\end{itemize}
\end{block}

\begin{block}{PHP}
\begin{itemize}
\item Powers more than 75\% of the Web
\item Had nothing for automated unit test generation
\item Had no types
\item Is interpreted, sources are always available
\end{itemize}
\end{block}

\end{frame}

\begin{frame}
\frametitle{Contract-Driven Testing}

\begin{block}{Principle}
Exploits contracts for test purposes:
\begin{itemize}
\item uses preconditions to generate test data
\item uses postconditions and invariants to establish test verdict by runtime
assertion
checking
\end{itemize}
\end{block}

\begin{block}{Contracts}
\begin{itemize}
\item Invented by B. Meyer in 1992 with Eiffel language
\item Describe a model using annotations
\item Express formal constraints: pre-, postconditions and invariants
\item Can be included directly in the source code applied to:
  \begin{itemize}
  \item classes attributes
  \item methods arguments
  \end{itemize}
\end{itemize}
\end{block}

\end{frame}

\begin{frame}
\frametitle{Design-by-Contract}

\begin{block}{Semantics of contracts}
\begin{itemize}
\item Contractual agreement:
  \begin{itemize}
  \item caller commits to satisfy the pre-condition
  \item called commits to establish its post-condition
  \end{itemize}
\item Invariants must be satisfied before and after the execution of the methods
\end{itemize}
\end{block}


\begin{block}{Issue of contracts}
\begin{itemize}
\item often expressed with logic formulæ
\item hard to generate data
\end{itemize}
\end{block}

\end{frame}

\begin{frame}
\frametitle{Previous works}

\begin{itemize}
\item Realistic domains
  \begin{itemize}
  \item structures to automate the validation and the generation of test data
  \end{itemize}
\item Praspel, a new specification language
  \begin{itemize}
  \item adopts Design-by-Contract paradigm
  \item based on realistic domains
  \item implementation in PHP for PHP
  \end{itemize}
\item Automated unit test generator: Praspel tool
  \begin{itemize}
  \item uses Praspel to perform Contract-Driven Testing
  \end{itemize}
\end{itemize}

\end{frame}

\begin{frame}
\frametitle{Motivations and contributions}

\begin{block}{Representing complex textual data}
\begin{itemize}
\item Regular expressions are not powerful enough (regular language)
\item We use grammar to represent these data (algebraic language)
\end{itemize}
\end{block}

\begin{block}{Contributions}
\begin{itemize}
\item Introduction of grammar-based testing in the Praspel tool
  \begin{itemize}
  \item generate and validate complex textual data
  \item PP, a new grammar description language
  \end{itemize}
\item New realistic domains: \code{grammar()} and \code{regex()}
\end{itemize}
\end{block}

\end{frame}

\outline

\section{Context}

\outlinereminder

\subsection{Realistic domains for PHP}

\begin{frame}
\frametitle{About realistic domains}

\begin{block}{Definition and goal}
\begin{itemize}
\item Are intended to be used for test generation purposes
\item Specify a set of relevant values that can be assigned to a data for a
specific context in a given program
\item Provide features for the validation and generation of data values
\end{itemize}
\end{block}

\begin{block}{Two important features}
\begin{itemize}
\item \textbf{Predicability}, checks if a value belongs to the realistic domain
\item \textbf{Samplability}, generates values that belong to the realistic
domain
\end{itemize}
\end{block}
The sampler can be of many kinds: a random generator, an iterator…
Features are user-defined

\end{frame}

\begin{frame}[fragile]
\frametitle{Realistic domains in PHP}

\begin{block}{Implementation}
Realistic domains as classes providing at least two methods:
\begin{itemize}
\item \code{predicate(\$q)}, takes a value \code{\$q} as input, returns a
boolean indicating the membership of the value to the realistic domain
\item \code{sample(\$sampler)}, generates values that belong to the realistic
domain according to a basic numeric-sampler \code{\$sampler}
\end{itemize}
\end{block}

\begin{block}{Hierarchical inheritance}
PHP realistic domains can inherit from each other, thanks to the PHP object
programming paradigm
\begin{center}
\includegraphics[width=6.5cm]{Figures/Realdom_univers.pdf}
\end{center}
\end{block}

\end{frame}

\begin{frame}[fragile,t]
\frametitle{User-defined realistic domain}

\begin{columns}[t]

\begin{column}{.5\textwidth}
\begin{block}{How to write your own realistic domain?}
\begin{itemize}
\item May extend an existing realistic domain
\item Write the \code{predicate(\$q)} method to add constraint on the
data sampled by the parent sampler
\item Write a new \code{sample(\$sampler)} method (optional)
\end{itemize}
\end{block}
\end{column}

\begin{column}{.5\textwidth}
\begin{exampleblock}{User-defined realistic domain}
{
\scriptsize
\begin{verbatim}
class Email extends String {

  public function predicate ( $q ) {

    $emailRegex = '...';
    return parent::predicate$(q)
        && 0 !== preg_match($emailRegex, $q);
  }

  public function sample ( $sampler ) {

    $characters  = array(...);
    $domains     = array(...);
    return $data = ...;
  }
}
\end{verbatim}
}
\code{email()} is intended to contain all email addresses
\end{exampleblock}
\end{column}

\end{columns}

\end{frame}

\begin{frame}
\frametitle{Parameters}

\begin{block}{Principle}
Realistic domains can receive parameters of many kinds: constants or realistic
domains themselves
\end{block}

\begin{exampleblock}{Constant arguments and realistic domains as arguments}
\begin{itemize}
\item \code{boundinteger(7, 42)} contains all the integers between 7 and 42
\item \code{string(boundinteger(4, 12), 0x20, 0x7e)} is intended to contain all
the strings of length between 4 and 12 constitued of characters from \code{0x20}
to \code{0x7e} (Unicode code-points)
\end{itemize}
\end{exampleblock}

\end{frame}

\subsection{Implementation in Praspel}

\begin{frame}
\frametitle{Presentation of Praspel}

\begin{block}{Praspel = PHP Realistic Annotation and SPEcification Language}
\begin{itemize}
\item Written in the API documentation (\code{/** ... */}) of the PHP code
\item Expresses contracts using formal constraints, called clauses, like:
  \begin{itemize}
  \item \code{@invariant}, class invariant on class attributes
  \item \code{@requires}, method precondition on class attributes and method
  arguments
  \item \code{@ensures}, method postcondition on class attributes, and method
  arguments and result
  \item \code{@throwable}, list of throwable exceptions by the method
  \end{itemize}
\end{itemize}
\end{block}

\begin{block}{Language properties}
\begin{itemize}
\item Assignment of realistic domains to a given data (\code{:})
\item A predicate $\backslash$\code{pred(…) }(expressed in the PHP
syntax), enriched with the $\backslash$\code{result} and
$\backslash$\code{old($e$)} constructs
\end{itemize}
\end{block}

\end{frame}

\begin{frame}[fragile]
\frametitle{Class with annotations}

\begin{exampleblock}{Generic example}
{
\small
\begin{verbatim}
class C {

    /**
     * @invariant _foo: float();
     */
    protected $_foo = 0;

    /**
     * @requires  baz:     … or … and
     *            qux:     … or … or …;
     * @ensures   \result: …;
     * @throwable AnException, AnotherException;
     */
     public function bar ( $baz, $qux ) {

        return …;
     }
}
\end{verbatim}
}
\end{exampleblock}

\end{frame}

\begin{frame}[fragile]
\frametitle{Praspel clauses}

\begin{exampleblock}{Example of a short Praspel contract}
{
\small
\begin{verbatim}
/**
 * @requires needle:   integer() and
 *           haystack: array([to integer()], boundinteger(1, 256));
 * @ensures  \result:  boolean();
 */
public function exists ( $needle,  $haystack ) {

    $intersect = array_intersect($haystack, array($needle));

    return 0 < count($intersect);
}
\end{verbatim}
}
\end{exampleblock}

\end{frame}

\subsection{Automated unit test generator}

\begin{frame}
\frametitle{Unit test generator and test verdict}

\begin{block}{Contract-Driven Testing}
The testing process works with the two features provided by the realistic
domains:
\begin{itemize}
\item random test data generation uses the sampler of each realistic domain
composing the precondition in order to satisfy it ({\em samplability})
\item test verdict is given by calling the predicate of each realistic domain
composing the postcondition ({\em predicability})
\end{itemize}
\end{block}

\begin{block}{Runtime Assertion Checking and test verdict}
\begin{itemize}
\item The RAC is performed by instrumenting the initial PHP code with additional
code that checks the contract clauses
\item Detected failures can be of five kinds: precondition, postcondition,
throwable, invariant or internal precondition (propagation) failure
\end{itemize}
\end{block}

\end{frame}

%\begin{frame}[fragile]
%\frametitle{Praspel clauses}
%
%\begin{exampleblock}{Example of a short Praspel contract}
%{
%\small
%\begin{verbatim}
%/**
% * @requires needle:   integer() and
% *           haystack: array([to integer()], boundinteger(1, 256));
% * @ensures  \result:  boolean();
% */
%public function exists ( $needle,  $haystack ) {
%
%    $intersect = array_intersect($haystack, array($needle));
%
%    return 0 < count($intersect);
%}
%\end{verbatim}
%}
%\begin{overprint}
%\onslide<1>
%foo
%
%\onslide<2>
%bar
%\end{overprint}
%\end{exampleblock}
%
%\end{frame}

\begin{frame}
\frametitle{Implementation in the Praspel tool}

\begin{block}{Environment for unit testing}
\begin{itemize}
\item Extensible and modular framework for generating and executing online tests
with a random data generator and runtime assertion checker
\item Praspel and its tools are freely available in Hoa
(\url{http://hoa-project.net}), a set of libraries for PHP
\end{itemize}
\end{block}
\begin{center}
\includegraphics[width=9cm]{Figures/Praspel_execution.png}
\end{center}

\end{frame}

\section{Grammar-based Testing}

\outlinereminder

\subsection{A new grammar description language}

\begin{frame}
\frametitle{Features of the PP language}

We propose the PP ({\em PHP Parser}) language as a new grammar description
language because none exists before in PHP

\begin{block}{Token}
\begin{itemize}
\item Lexical unit
\item Represented by the PCRE ({\em Perl Compatible Regular Expression})
\item Namespaces (operator \code{->} to change namespace)
\end{itemize}
\end{block}

\begin{block}{Rule}
\begin{itemize}
\item Identified by a name
\item Sequence of tokens is based on the following operators:
  \begin{itemize}
  \item repetition: $e$\code{\{}$x$\code{,}$y$\code{\}}, $e$\code{?},
  $e$\code{+}, $e$\code{*}
  \item concatenation: $e_1 \dots e_i \dots e_n$
  \item disjunction and grouping: $e_1$ \code{$\vert$} $\dots$ \code{$\vert$}
  $e_i$ \code{$\vert$} $\dots$ \code{$\vert$} $e_n$ and \code{(}$e$\code{)}
  \item token: \code{<}$t$\code{>} or \code{::}$t$\code{::}
  \item call a rule: $r$\code{()}
  \item add a marker: \code{\#}$n$
  \end{itemize}
\end{itemize}
\end{block}

\end{frame}

\begin{frame}[fragile]
\frametitle{Syntax}

\begin{exampleblock}{Simplified XML grammar expressed with PP}
{
\footnotesize
\begin{verbatim}
%skip   s          \s
%token  lt         <      -> tag
%token  tag:skip   \s
%token  tag:name   \w+
%token  tag:slash  /
%token  tag:gt     >      -> default
%token  content    [^<]+

xml:
    ::lt:: <name>
  ( ::slash:: ::gt:: #fold
  | ::gt:: ( <content> | xml()+ )? ::lt:: ::slash:: ::name:: ::gt:: #unfold )
\end{verbatim}
}
Valid: \code{<a> <b>foo</b> <b>bar</b> <c /> </a>}, \\
but also: \code{<a>foo</z>}
\end{exampleblock}

\end{frame}

\begin{frame}[fragile]
\frametitle{Unification}

\begin{block}{Principle}
Unification expresses another constraint in grammar. \\
All tokens $t$\code{[}$i$\code{]} with the same $i$ have the same value locally
to a rule instance
\end{block}

\begin{exampleblock}{Unified XML tag names}
{
\footnotesize
\code{xml:} \\
\code{\quad\quad::lt:: <name\red{[0]}>} \\
\code{\quad( ::slash:: ::gt:: \#fold} \\
\code{\quad| ::gt:: ( <content> | xml()+ )? ::lt:: ::slash:: ::name\red{[0]}:: ::gt:: \#unfold )} \\
}
Invalid: \code{<a>foo</z>}
\end{exampleblock}

\end{frame}

\subsection{Grammar-based realistic domain}

\begin{frame}
\frametitle{Grammar-based realistic domain}

Such a realistic domain has also two features:
\begin{itemize}
\item[\Square] \textbf{Predicability}, checks the conformance between the
data and the grammar
\item[\Square] \textbf{Samplability}, will generate a data matching the grammar
\end{itemize}

\end{frame}

\begin{frame}
\frametitle{Predicability process}

\begin{center}
\includegraphics[width=8cm]{Figures/Compiling_process.pdf}
\end{center}
\begin{itemize}
\item Lexing: transform a given data into a sequence of tokens
\item Parsing: analyze this sequence according to the rules
  \begin{itemize}
  \item derive from left to right and top to bottom
  \item grammar is ambiguous: $LL(*)$, implies backtracks
  \end{itemize}
\end{itemize}

\end{frame}

\begin{frame}
\frametitle{Abstract Syntax Tree}

Classic compilation process ends by building an AST, which accepts visitors
(design-pattern) \\
\begin{exampleblock}{\code{<a> <b>foo</b> <b>bar</b> <c /> </a>} associated AST}
\begin{center}
\includegraphics[width=6cm]{Figures/AST.pdf}
\end{center}
\end{exampleblock}
Useful for additional verifications that can not be expressed in the grammar

\end{frame}

\begin{frame}
\frametitle{Grammar-based realistic domain}

Such a realistic domain has also two features:
\begin{itemize}
\item[\checked] \textbf{Predicability}, checks the conformance between the
data and the grammar
  \begin{itemize}
  \item ensured by the parsing process
  \end{itemize}
\item[\Square] \textbf{Samplability}, will generate a data matching the grammar
  \begin{enumerate}
  \item generate tokens
  \item generate sequences of tokens
  \end{enumerate}
\end{itemize}

\end{frame}

\begin{frame}[t]
\frametitle{Generate tokens}

\begin{columns}[t]

\begin{column}{.55\textwidth}
\begin{exampleblock}{Example}
\onslide<+->
\code{\lighter{([ae]+|[x-z]!)}\{1,3\}} \\
\onslide<+->
\quad $\rightarrow$ \code{([ae]+|[x-z]!)\lighter{([ae]+|[x-z]!)}} \\
\onslide<+->
\quad $\rightarrow$ \code{([ae]+)\lighter{([ae]+|[x-z]!)}} \\
\onslide<+->
\quad $\rightarrow$ \code{[ae][ae]\lighter{([ae]+|[x-z]!)}} \\
\onslide<+->
\quad $\rightarrow$ \code{e[ae]\lighter{([ae]+|[x-z]!)}} \\
\onslide<+->
\quad $\rightarrow$ \code{ea\lighter{([ae]+|[x-z]!)}} \\
\onslide<+->
\quad $\rightarrow$ \code{\lighter{ea}([x-z]!)} \\
\onslide<+>
\quad $\rightarrow$ \code{eay!} \\
\end{exampleblock}
\end{column}

\begin{column}{.45\textwidth}
\onslide<1->
\begin{block}{Approach}
\begin{itemize}
\item Random and uniform choices
\item Isotropic exploration
\item Naive but tokens are not important here, sequences of tokens are important
\end{itemize}
\end{block}

\begin{block}{Repetition unfolding}
Upper bound of \code{+} and \code{*} is set to a predefined integer
\end{block}
\end{column}

\end{columns}

\end{frame}

\begin{frame}
\frametitle{Rules to generate sequences of tokens}

We propose 3 differents algorithms

\begin{block}{Why?}
\begin{itemize}
\item because we can (thanks to Praspel)
\item because an algorithm to rule them all does not exist
  \begin{itemize}
  \item to each context of use is associated an algorithm
  \item we retain 3 algorithms from the literature
  \end{itemize}
\end{itemize}
\end{block}

\begin{block}{Algorithms}
\begin{itemize}
\item Random and uniform generation
\item Bounded exhaustive generation
\item Coverage-based generation
\end{itemize}
\end{block}

\end{frame}

\begin{frame}[t]
\frametitle{Random and uniform generation}

\vspace*{-.7cm}

\begin{columns}[t]

\begin{column}{.6\textwidth}
\begin{exampleblock}{Example: \code{f()}, \code{g()} and $n = 5$}
\code{f: <a> g()} \\
\code{g: ( <b> <c> | <d> | f() )\{1,3\}} \\
\vspace{2em}

\begin{overprint}

\onslide<1>
\scriptsize
\begin{tabular}{rcl}
$f(5)$ & = & $1 . g(4)$ \\
$g(4)$ & = & $(0 + 0 + f(4)$ \\
& & $+ ( 0 + 0 + f(3) ) . ( 0 + 1 + f(1) )$ \\
& & $+ ( 1 + 0 + f(2) ) . ( 1 + 0 + f(2) )$ \\
& & $+ ( 0 + 1 + f(1) ) . ( 0 + 0 + f(3) )$ \\
& & $+ ( 1 + 0 + f(2) ) . ( 0 + 1 + f(1) ) . ( 0 + 1 + f(1) )$ \\
& & $+ ( 0 + 1 + f(1) ) . ( 1 + 0 + f(2) ) . ( 0 + 1 + f(1) )$ \\
& & $+ ( 0 + 1 + f(1) ) . ( 0 + 1 + f(1) ) . ( 1 + 0 + f(2) )$ \\
$f(4)$ & = & $1 . g(3)$ \\
$g(3)$ & = & $…$
\end{tabular}
%p(5, f) = 1 . g(4) \\
%        = 20
%p(4, g) =   ( 0 + 0 + p(4, f) ) \\
%    3     + ( 0 + 0 + p(3, f) ) . ( 0 + 1 + p(1, f) ) \\
%    4     + ( 1 + 0 + p(2, f) ) . ( 1 + 0 + p(2, f) ) \\
%    3     + ( 0 + 1 + p(1, f) ) . ( 0 + 0 + p(3, f) ) \\
%    2     + ( 1 + 0 + p(2, f) ) . ( 0 + 1 + p(1, f) ) . ( 0 + 1 + p(1, f) ) \\
%    4     + ( 0 + 1 + p(1, f) ) . ( 1 + 0 + p(2, f) ) . ( 0 + 1 + p(1, f) ) \\
%    1     + ( 0 + 1 + p(1, f) ) . ( 0 + 1 + p(1, f) ) . ( 1 + 0 + p(2, f) ) \\
%        = 3 + 3 + 4 + 3 + 2 + 4 + 1 \\
%        = 20 \\
%
%p(4, f ) = 1 . g(3) \\
%         = 3 \\
%
%p(3, g) =   ( 0 + 0 + p(3, f) ) \\
%    0     + ( 0 + 0 + p(2, f) ) . ( 0 + 0 + p(1, f) ) \\
%    0     + ( 0 + 0 + p(1, f) ) . ( 0 + 0 + p(2, f) ) \\
%        = 3 \\
%
%p(3, f) = 1 . g(2) \\
%        = 3 \\
%
%p(2, g) =     1 + 0 + p(2, f) \\
%          + ( 0 + 1 + p(1, f) ) . ( 0 + 1 + p(1, f) ) \\
%        = 3 \\
%
%p(2, f) = 1 . g(1) \\
%        = 1 \\
%p(1, g) = 1 \\
%p(1, f) = 0

\onslide<2->
\footnotesize
example:
\begin{center}
\begin{tabular}{r|c|c|c|c|c|}
         &  5 &  4 & 3 & 2 & 1 \\
\hline
\code{f} & 24 &  3 & 3 & 1 & 0 \\
\hline
\code{g} &  - & 24 & 3 & 3 & 1 \\
\hline
\end{tabular}
\end{center}

choice-point and probability: \\
\code{h(3) = 20} \\
\code{i(2) = 6} \\
\code{j(2) = 14} \\
\begin{center}
\code{h: <x> ( }
$\underbrace{\code{i()}}_{\dfrac{6}{20}}$\code{ | }
$\underbrace{\code{j()}}_{\dfrac{14}{20}}$\code{ )}
\end{center}
\end{overprint}

\end{exampleblock}
\end{column}

\begin{column}{.4\textwidth}
\begin{block}{Approach}
\begin{itemize}
\item An expected sequence size $n$ and uniform probability distribution among
all the possible sequences
\item Recursive method to count all possible sub-structures of size $n$
\item Counting helps to compute cumulative distribution functions, which
guide exploration
\end{itemize}
\end{block}

\begin{block}{Repetition unfolding}
Upper bound of \code{+} and \code{*} is set to $n$
\end{block}
\end{column}

\end{columns}

\end{frame}

\begin{frame}
\frametitle{Random and uniform generation}

\begin{block}{Counting function $\psi$}
\begin{eqnarray*}
\psi(n, e) & = & \delta_n^1 \text{ \ \ if $e$ is a token} \\
\psi(n, e_1 \;\cdot\; \dots \;\cdot\; e_k) & = & \sum_{\gamma \,\in\, \Gamma_k^n}\;\;
                                                 \prod_{\alpha = 1}^k\;\;
                                                 \psi(\gamma_\alpha, e_\alpha) \\
\psi(n, e_1 \;\vert\; \dots \;\vert\; e_k) & = & \sum_{\alpha\,=\,1}^k
                                                 \psi(n, e_\alpha) \\
\psi(n, e^{\{x,y\}}) & = & \sum_{\alpha\,=\,x}^y\;\;
                           \sum_{\gamma \,\in\, \Gamma_\alpha^n}\;\;
                           \prod_{\beta\,=\,1}^\alpha \psi(\gamma_\beta, e) \\
                     &   & \text{with $0 \le x \le y$}
\end{eqnarray*}
\end{block}

\end{frame}

\begin{frame}[t]
\frametitle{Bounded exhaustive generation}

\begin{columns}[t]

\begin{column}{.55\textwidth}
\begin{exampleblock}{Example: \code{f()}, \code{g()} and $n = 10$}
\code{f: <a> g()} \\
\code{g: ( <b> <c> | <d> | f() )\{1,3\}} \\

\begin{enumerate}
\onslide<+->
\item \code{<a> <b> <c>}
\onslide<+->
\item \code{<a> <d>}
\onslide<+->
\item \code{<a> <a> <b> <c>}
\onslide<+->
\item \code{<a> <a> <d>}
\onslide<+->
\item \code{<a> <a> <a> <b> <c>}
\onslide<+->
\item \code{<a> <a> <a> <d>}
\onslide<+->
\item \code{<a> <a> <a> <a> <a> <b> <c>}
\item …
\end{enumerate}
\onslide<+->
more than 10000 solutions
\end{exampleblock}
\end{column}

\onslide<1->
\begin{column}{.45\textwidth}
\begin{block}{Approach}
\begin{itemize}
\item Generate all possible sequences (exhaustive) up to a given size $n$
(bounded)
\item The algorithm behaves like an iterator
\item Based on multiset (set with repetition)
\end{itemize}
\end{block}

\begin{block}{Repetition unfolding}
Upper bound of \code{+} and \code{*} is set to $n$
\end{block}
\end{column}

\end{columns}

\end{frame}

\begin{frame}
\frametitle{Bounded exhaustive generation}

\begin{block}{Function $\beta$ for bounded exaustive generation}
\footnotesize
\begin{eqnarray}
\beta(1, e) & = & \{\code{sample}(e)\} \text{ \ \ if $e$ is a token} \nonumber \\ 
\beta(n, e) & = & \{\} \text{ \ \ if } n \neq 1 \nonumber \\
\beta(n, e_1 \;\vert\; e_2) & = & \beta(n,e_1) \cup \beta(n,e_2) \nonumber \\
\beta(n, e_1 \;\cdot\; e_2) & = & \bigcup_{p\,=\,1}^{n-1}
                                  \beta(p,e_1) \cdot \beta(n - p,e_2) \nonumber \\
\beta(n, e^{\{x,y\}}) & = & \bigcup_{p\,=\,x}^{y} \beta(n,e^{p}) \nonumber \\
\beta(n, e^{\star}) & = & \bigcup_{p\,=\,0}^{n} \beta(n,e^{p}) \nonumber \\
\beta(n,e^{+}) & = & \beta(n, e \;\cdot\; e^{\star}) \nonumber \\
\beta(n,e^{0}) & = & \{\} \nonumber \\
\beta(n,e^{1}) & = & \beta(n,e) \nonumber \\
\beta(n,e^{p}) & = & \beta(n, e \;\cdot\; e^{p-1})\text{ \ \ if } p \ge 2
\nonumber
\end{eqnarray}
\end{block}

\end{frame}

\begin{frame}[t]
\frametitle{Coverage-based generation}

\begin{columns}[t]

\begin{column}{.4\textwidth}
\begin{exampleblock}{Example: \code{f()}, \code{g()}}
\small
\code{f: <a> g()} \\
\code{g: ( <b> <c> | <d> | f() )\{1,3\}} \\

\begin{enumerate}
\onslide<+->
\item \code{<a> <d> <b> <c> <a> <d>}
\end{enumerate}
1 solution
\end{exampleblock}

\onslide<1->
\begin{block}{Repetition unfolding}
\begin{itemize}
\item \code{*} is bounded to 0, 1 or 2
\item \code{+} is unfolded 1 or 2 times
\item \code{\{x,y\}} is unfolded $x$, $x + 1$, $y - 1$ and $y$ times
\end{itemize}
\end{block}
\end{column}

\onslide<1->
\begin{column}{.6\textwidth}
\begin{block}{Approach}
\begin{itemize}
\item Reduce the combinatorial explosion
\item Aims at producing data that activate all the branches of the grammar rules
\item A rule is said to be covered if and only if its sub-rules have all been
covered
\item A token is said to be covered if it has been successfully used in a data
generation
\item To ensure diversity, a random choice is made amongst the remaining
sub-rules of a choice-point to cover
\item Boundary test generation heuristics to avoid combinatorial explosion and
guarantee the termination
\end{itemize}
\end{block}
\end{column}

\end{columns}

\end{frame}

\begin{frame}
\frametitle{Coverage-based generation}

\begin{block}{Coverage-based function $\phi$}
\begin{eqnarray*}
\phi(p, e) & = & [\,\code{sample}(e)\,] \text{ \ \ when $e$ is a token} \\
\phi(p, e_1 \;\cdot\; e_2) & = & \phi(\phi(p, e_1),e_2) \\
\phi(p, e_1 \;\vert\; \ldots \;\vert\; e_k) & = & \phi(p,e_1) \;\oplus\; \ldots \;\oplus\; \phi(p,e_k) \\
\phi(p, e^?) & = & [\,] \;\oplus\; \phi(p,e) \\
\phi(p, e^\star) & = & [\,] \;\oplus\; \bigoplus_{i\,=\,1}^{\infty} \phi(p, \underbrace{e \;\cdot\; \dots \;\cdot\; e}_i)  \\
\phi(p, e^+) & = & \bigoplus_{i\,=\,1}^{\infty} \phi(p, \underbrace{e \;\cdot\; \dots \;\cdot\; e}_i)  \\
\phi(p, e^{\{x,y\}}) & = & \bigoplus_{i\,=\,x}^y \phi(p, \underbrace{e \;\cdot\; \dots \;\cdot\; e}_i)
\end{eqnarray*}
\end{block}

\end{frame}

\begin{frame}
\frametitle{Pros and cons}

\begin{block}{Complex textual data generation}
\begin{itemize}
\item Random and uniform generation:
  \begin{itemize}
  \item[\green \smiley] fast for small data, diversity of data and fixed size
  \item[\frownie] counting phase is exponential ($n > 10$ needs at least 2 hours),
  despite that generation is fast
  \end{itemize}
\item Bounded exhaustive generation:
  \begin{itemize}
  \item[\green \smiley] fast for small data and exhaustiveness is efficient
  \item[\frownie] exponential number of data
  \end{itemize}
\item Coverage-based generation:
  \begin{itemize}
  \item[\green \smiley] fast for medium and big data and diversity of data
  \item[\frownie] do not consider size of data
  \end{itemize}
\end{itemize}
\end{block}

\end{frame}

\begin{frame}
\frametitle{Grammar as a realistic domain}

Such a realistic domain has also two features:
\begin{itemize}
\item[\checked] \textbf{Predicability}, checks the conformance between the
data and the grammar
  \begin{itemize}
  \item ensured by the parsing process
  \end{itemize}
\item[\checked] \textbf{Samplability}, will generate a data matching the grammar
  \begin{itemize}
  \item ensured by one of the three algorithms
  \end{itemize}
\end{itemize}

\end{frame}

\subsection{Experimentation}

\begin{frame}
\frametitle{Experimentation}

\begin{block}{Self-validation}
\begin{itemize}
\item Generate and validate data with PP and other parsers
\item Considered grammars: JSON and PCRE, with other parsers from Mozilla Gecko
and PHP
\item All produced data were parsed correctly
\end{itemize}
\end{block}

\begin{block}{Mutation}
\begin{itemize}
\item Then, we consider simple grammar mutation operators
\item Some generated data were parsed correctly by PP but not by
other parsers (due to the backtracking)
\item After fixing the bug, we performed the same kinds of experiments with PP
and other parsers
\end{itemize}
\end{block}

\end{frame}

\section{Conclusion}

\outlinereminder

\begin{frame}
\frametitle{Conclusion}

\begin{block}{What have we seen?}
\begin{itemize}
\item Realistic domains specifying data and providing two useful features for
automated test generation: predicability and samplability
\item Praspel, a new Design-by-Contract language implementing realistic domains
\item Praspel tool: a test generation and execution framework to automate unit
testing in PHP
\item Grammar-based Testing is introduced in Praspel
\item PP, a new grammar description language
\item Two new realistic domains join the standard library
\end{itemize}
\end{block}

\end{frame}

\begin{frame}
\frametitle{Conclusion}

\begin{block}{Future works}
\begin{itemize}
\item Extend case studies in order to evaluate the relevance of the
coverage-based test generation, in terms of fault detection and code coverage of
the system under test
\item Improve the generation algorithms so as to avoid rejection as much as
possible (look at UDITA)
\item Implement Praspel into other languages (e.g. Java, C, Javascript)
\end{itemize}
\end{block}

\end{frame}

\begin{frame}
\frametitle{Thanks!}

\begin{center}
Thank you for your attention! Any questions?
\end{center}

\end{frame}

\end{document}
